%% Vol I, Chapter 1 — Introduction and Architecture Overview
%% SOURCE: docs/working/architecture/03-architecture/OVERVIEW.md
%%         .claude/CLAUDE.md (Project Overview section)
%%         ONTOLOGY.md (thermodynamic lexicon)
%%         docs/working/papers/EXECUTIVE_SUMMARY.md

\chapter{Introduction}
\label{vol1:chap:intro}

%% TODO(bob): write the motivating paragraph — why a physics-grounded adaptive
%% runtime? What problem does this solve for StarshipOS?

\section{What StarForth Is}
\label{vol1:sec:intro:what}

%% TODO(bob): 2–3 paragraphs from EXECUTIVE_SUMMARY.md, calibrated for
%% a technical but non-FORTH audience. Key facts to include:
%%   - FORTH-79 compliant VM, strict ANSI C99, no GNU extensions
%%   - Primary execution engine for StarshipOS
%%   - Patent-pending adaptive runtime (cite james:2025:patent)
%%   - Published SSRN results (cite james:2025:ssrn)

\section{Technology Stack}
\label{vol1:sec:intro:stack}

%% TODO(bob): reproduce the three-layer diagram from OVERVIEW.md in TikZ or as
%% a table:
%%   StarshipOS (future) | LithosAnanke v1.0.8 | StarForth v3.0.3 | Linux/L4Re/bare metal

\begin{table}[ht]
  \centering
  \caption{StarshipOS technology stack}
  \label{tab:stack}
  \begin{tabular}{ll}
    \toprule
    Layer & Component \\
    \midrule
    Future OS       & StarshipOS (self-hosting) \\
    Kernel          & LithosAnanke v1.0.8 / StarKernel (bare-metal UEFI) \\
    VM              & StarForth v3.0.3 (FORTH-79, physics-driven adaptive RT) \\
    Host platform   & Linux, L4Re/Fiasco.OC, bare metal (amd64, aarch64, riscv64) \\
    \bottomrule
  \end{tabular}
\end{table}

\section{Thermodynamic Modeling Language}
\label{vol1:sec:intro:ontology}

%% TODO(bob): one paragraph explaining that the runtime uses thermodynamic
%% concepts as a modeling language, not as physics claims. Cite ONTOLOGY.md.
%% The key calibration sentence is in ACADEMIC_WORDING_GUIDELINES.md.

\section{Document Scope and Organization}
\label{vol1:sec:intro:scope}

%% TODO(bob): paragraph describing what each chapter covers (interpreter,
%% physics runtime, formal verification, experimental results, build, ACL).
